\subsection{14. Nachprüfung der beiden zuletzt zugeschickten
Perspektiven}\label{nachpruxfcfung-der-beiden-zuletzt-zugeschickten-perspektiven}

Die beiden Texte mit den Anfängen \textbf{„Ja. Nach dem erneuten
Querschnitt \ldots``} und \textbf{„Ich sehe keine bereits versteckte
Formel \ldots``} wurden vollständig gelesen. Ihre Vorschläge sind
Prüfhypothesen, keine übernommenen Arbeitsanweisungen oder bereits
bewiesene Resultate. Beide Originale liegen mit Hashes unverändert im
Prüfpaket. Die folgende Prüfung ergänzt die Spin-Lift-Rechnung, ersetzt
sie nicht.

\subsubsection{14.1 Tragender gemeinsamer Kern, aber keine automatische
Quellenauswahl}\label{tragender-gemeinsamer-kern-aber-keine-automatische-quellenauswahl}

Der positive Prozesskern

\[\adjustbox{max width=\linewidth}{$\displaystyle \Gamma(u,v)=\omega(u^\dagger v)$}\]

ist eine gute gemeinsame Beschreibung der erreichbaren Vorgänge. Ist
eine unital vorgegebene C*-Algebra mit ihrer Multiplikation und einem
Zustand gegeben, bestimmt die GNS-Konstruktion die zyklische Darstellung
bis auf eine zustands- und algebraerhaltende unitäre Äquivalenz. Bei
vollständig zeit- und instrumentmarkierten Kernen werden die so
markierten Prozesse mit rekonstruiert. Das ist die präzise Bedeutung von
„gleicher vollständiger Prozess``.

\textbf{Zwei Einschränkungen sind entscheidend.} Erstens genügt die
Positivität einer beliebigen Matrix über Wortnamen nicht: Ein Nullwort
muss auch nach erlaubter Linksmultiplikation null bleiben. Die positive
Gram-Matrix diag(1,0,1) über den Wörtern \((1,a,a^2)\) verletzt dies,
denn \([a]=0\), aber \([a^2]\ne0\). Ohne Kompatibilität mit den
Algebrarelationen bekommt man einen Vektorraum von Antworten, noch keine
ausführbare Algebra.

Zweitens enthält ein statischer Kern keine bisher unmarkierte physische
Zeit. Auf \(\mathbb C^3\) besitzen

\[\adjustbox{max width=\linewidth}{$\displaystyle \Omega=(1,0,0)^T,\qquad H_1=\operatorname{diag}(0,1,2),\qquad H_2=\operatorname{diag}(0,1,3)$}\]

denselben eindeutigen Grundvektor und für sämtliche zeitunmarkierten
Matrixwörter denselben Kern. Die Anregungsenergien und zeitabhängigen
Antworten sind verschieden. Nimmt man Zeittranslationen schon in die
Wortmarkierungen auf, wird diese fehlende Information \textbf{mit
eingegeben}. Die Rekonstruktion ist dann korrekt, aber keine Auswahl
dieser Zeittranslationen aus den statischen Daten.

Die vorgeschlagene Selbstabschlussbedingung ist deshalb ein
Forschungsprogramm, noch keine definierte eindeutige Gleichung.
Festzulegen sind die realisierbaren Controller, ihre Ressourcen und der
Abschlussoperator; außerdem muss eine nichttriviale Lösung bewiesen
werden. Reine Positivität oder „kleinste konsistente Lösung`` kann auch
eine stationäre skalare Theorie zulassen. Ein Quotient nach interner
Ununterscheidbarkeit spart redundante Darstellungen; er macht fehlende
Eingriffe nicht ausführbar.

\subsubsection{14.2 Die modulare Zeitidee ist relevant, aber noch nicht
unser fehlender
Hamiltonoperator}\label{die-modulare-zeitidee-ist-relevant-aber-noch-nicht-unser-fehlender-hamiltonoperator}

Die vorgeschlagene Verbindung zu Connes und Rovelli ist sachlich
richtig: Ihre
\href{https://arxiv.org/abs/gr-qc/9406019}{Thermal-Time-Arbeit}
postuliert eine physische Interpretation des vom Zustand bestimmten
modularen Flusses. Das ist eine physische Hypothese auf Grundlage eines
mathematischen Satzes, kein allgemeiner Beweis, dass jeder modulare
Fluss eine physische Uhr ist.

Für eine volle endliche Matrixalgebra und einen treuen Zustand
\(\rho>0\) gilt

\[\adjustbox{max width=\linewidth}{$\displaystyle \sigma_t^\rho(A)=\rho^{it}A\rho^{-it}.
\qquad K=-\log\rho.$}\]

Mit der Konvention \(\alpha_s^H(A)=e^{isH}Ae^{-isH}\) gilt

\[\adjustbox{max width=\linewidth}{$\displaystyle \boxed{\sigma_t^\rho=\alpha_{-\beta t}^H\ \text{für alle }t
\iff K=\beta H+cI
\iff \rho=Z^{-1}e^{-\beta H}.}$}\]

Beweis: Ableiten bei null ergibt \([K-\beta H,A]=0\) für jede Matrix A.
Der Kommutant der vollen Matrixalgebra besteht aus Skalaren. Der Rückweg
folgt direkt durch Exponentiation. Für eine direkte Summe von
Matrixalgebren ersetzt ein zentrales Element die skalare Konstante.

Damit ist ein konkreter Prüfmaßstab vorhanden: Ein \textbf{unabhängig
aus der Quelle gewonnener} Zustand müsste diese Beziehung erfüllen. Wird
zunächst aus H ein Gibbszustand hergestellt und danach \(-\log\rho\)
berechnet, wurde H nicht hergeleitet. Selbst Stationarität genügt nicht:
Für \(H=\operatorname{diag}(0,1,2)\) ist
\(\rho=\operatorname{diag}(1/2,1/3,1/6)\) treu und kommutiert mit H; die
modularen benachbarten Frequenzen sind jedoch \(\log(3/2)\) und
\(\log 2\), während H gleiche Energieabstände hat. Keine gemeinsame
Zeitskalierung macht beide Flüsse gleich.

\textbf{Auf dem hier bewiesenen nativen Modellgrundzustand scheitert die
unmittelbar globale Variante bereits früher.} Ein reiner Grundzustand
ist auf der vollen Matrixalgebra seines Ladungssektors nicht treu: Ein
nichtnull Projektor orthogonal zu ihm hat Erwartungswert null.
Komprimiert man auf seinen Träger, bleibt nur eine eindimensionale
Algebra und trivialer modularer Fluss. Das ist keine Einschränkung der
Vakuum-Modulartheorie lokaler unendlicher Algebren. Ein reiner
Gesamtzustand kann auf einer echten Teilalgebra treu sein; das Beispiel

\[\adjustbox{max width=\linewidth}{$\displaystyle \Psi=\sqrt{2/3}|00\rangle+\sqrt{1/3}|11\rangle,
\qquad\rho_A=\operatorname{diag}(2/3,1/3)$}\]

zeigt genau diesen Unterschied. Die Wahl und physische Verfügbarkeit
dieser Teilalgebra müssen aber aus der Quelle folgen. Für den
algebraabhängigen Begriff von Verschränkung und zyklisch-separierenden
Vektoren siehe \href{https://arxiv.org/abs/1803.04993}{Wittens
Darstellung}.

Noch eine verbleibende Mehrdeutigkeit: Auf \(\mathbb C\oplus M_2\)
erzeugen die treuen Zustände

\[\adjustbox{max width=\linewidth}{$\displaystyle \rho_w=w\oplus(1-w)\rho_2,\qquad0<w<1$}\]

denselben modularen Fluss bei festem \(\rho_2\). Die Blockgewichte
verschwinden in der Konjugation. Der Fluss bewegt keine zentralen
Superselektionsgewichte und bestimmt sie nicht. Der maximal gemischte
Zustand ist zwar besonders symmetrisch und treu, liefert aber auf einer
Matrixalgebra überhaupt keinen nichttrivialen modularen Zeitfluss. „Nimm
den einfachsten Zustand`` allein löst daher unsere Quellen-/Uhrfrage
nicht.

\subsubsection{14.3 Neue Schranke für den vorgeschlagenen
Kausalitätstest}\label{neue-schranke-fuxfcr-den-vorgeschlagenen-kausalituxe4tstest}

Die Konstruktion von Raumzeit aus relativen modularen Flüssen ist eine
ernsthafte Richtung. Sie verlangt zusätzliche Relationen der
\textbf{Algebren}, nicht bloß einige ähnliche Eigenwerte. Ein relevanter
mathematischer Rahmen sind halbseitige modulare Inklusionen; einen
präzisierten Struktursatz geben
\href{https://arxiv.org/abs/math/0412061}{Araki und Zsido}.

Für die endlichen Testalgebren gilt jedoch der einfache allgemeine Satz:
Ist \(\mathcal N\) endlichdimensional und

\[\adjustbox{max width=\linewidth}{$\displaystyle \sigma_t(\mathcal N)\subseteq\mathcal N\qquad(t\ge0),$}\]

dann ist diese Inklusion für jedes t eine \textbf{Gleichheit}. Denn eine
Automorphie erhält die Vektorraumdimension. Inklusion plus gleiche
endliche Dimension erzwingt Gleichheit; die inverse Automorphie gibt die
Gleichheit auch für negative t. Eine echte einseitige Schrumpfung der
Algebra ist so unmöglich.

Dieser Dimensionsbeweis ist nicht ein globaler Ausschluss modularer
Raumzeitrekonstruktion. Er zeigt, warum ein einzelner endlicher
Clockblock kein nichttriviales Beispiel der benötigten Halbseitigkeit
liefern kann. Man braucht eine kontrollierte unendliche Algebra oder
einen Grenzübergang, in dem die endliche Dimensionsschranke nicht mehr
greift. Gerade diese Kontrolle darf nicht durch einen großen endlichen
Fit ersetzt werden.

\subsubsection{14.4 Phasen, Gesamtsymmetrie und Unsichtbarkeit kleiner
Tests}\label{phasen-gesamtsymmetrie-und-unsichtbarkeit-kleiner-tests}

Der zweite Text trägt in drei Punkten unmittelbar:

\begin{itemize}
\tightlist
\item
  \textbf{Die Schleifenphase ist entscheidende Information.} Die zwei
  Dreiecksmatrizen wurden erneut unabhängig exakt nachgerechnet:
  Fermionränge 3 und 2, Überlappungsprodukte \(+1/8\) und \(-1/8\), aber
  identische Paarmatrix mit Spektrum \((12,6,6)\). Abschnitte 3-7 gehen
  darüber hinaus: Für den Viererclock wird das verlorene Vorzeichen an
  höhere Antworten, einen Grundzustand und dessen geladene Linien
  angeschlossen. Eine bloße Übereinstimmung von Paarzahlen reicht
  tatsächlich nicht.
\item
  \textbf{Volle Quelle statt isolierter Tensor.} Für tatsächlich
  gesetzte Strukturen gilt
  \(\operatorname{Stab}(W,\text{Clock},\ldots)\subseteq\operatorname{Stab}(W)\).
  Eine reale Markierung kann die Gruppe einschränken; eine umbenannte
  Basis nicht. Der neue Viererclock wirkt nur im zusätzlichen Kopienraum
  und erhält gerade die volle innere Gruppe. Er löst daher keine unter
  innerer G-Invarianz bewiesene Kontrollsperre durch Symmetriebruch. Bei
  der ursprünglichen vollständigen Quelle bleibt das gemeinsame
  Stabilisatorproblem offen. Die bisherige Schur-Argumentation für
  unseren Kandidaten wird durch diese Prüfspur nicht entkräftet.
\item
  \textbf{Endlich viele Ladungssektoren wählen kein unbegrenztes Gesetz
  aus.} Wenn alle Versuche \(N\le N_{\max}\) erhalten, setze
  \(m=\lfloor N_{\max}/2\rfloor+1\). Dann ist
  \(F_m(N_b)=\prod_{j=0}^{m-1}(N_b-j)\) dort identisch null, auf
  ganzzahligen Bosonbesetzungen überall nichtnegativ und erstmals bei
  \(N_b=m\) ungleich null. Folglich ist \(H+\eta F_m(N_b)\) bei
  \(\eta>0\) in sämtlichen getesteten Abläufen ununterscheidbar von H,
  solange die Instrumente diese Sektoren nicht verlassen. Die
  Deformation respektiert Gesamtladung und innere Symmetrie. Der
  allgemeine Beweis ist das Nullprodukt; die Reproduktion ergänzt exakte
  Beispiele bis \(N_{\max}=12\). Ein begründeter Grad-/Operationsvertrag
  kann die Deformation ausschließen, kleine Tests allein nicht.
\end{itemize}

Die Universalitätsalternative aus beiden Texten bleibt sinnvoll: Für
makroskopische Gesetze muss möglicherweise eine robuste Klasse und nicht
genau ein Mikrograph ausgewählt werden. Das ersetzt die
RG-/Kontinuumsprüfung nicht. Ein Code ist außerdem ein möglicher
Kommunikationstest, kein universell vorgeschriebenes Naturgesetz.

\subsubsection{14.5 Konsequenz für die weitere
Arbeit}\label{konsequenz-fuxfcr-die-weitere-arbeit}

Die Priorität wird präziser, nicht breiter: \textbf{Quelle samt
fermionischem Lift und realem Operationssatz bestimmen; auf ihr Zustand
und Antworten gemeinsam berechnen; dann lokale Algebren und ihre
Skalierung prüfen.} Ein unabhängig gewonnener treuer Teilzustand und
eine nichttriviale modulare Vergleichsrelation wären ein echter
zusätzlicher Anschluss. Ein aus H definierter Gibbszustand oder ein
bereits zeitmarkierter vollständiger Prozesskern wäre dagegen eine
konsistente Umformulierung, noch keine neue Herleitung der Dynamik.

Keiner der beiden Texte liefert bereits die fehlenden Voraussetzungen
für eine vollständige TOE, RH, effiziente Faktorisierung oder P versus
NP. Ihre tragenden Hinweise sind integriert, ihre offenen Auswahlfragen
bleiben sichtbar. Die neuen kleinen Gegenrechnungen sind in
\codeword{verify\_new\_aspects.py} von der großen
Spin-Lift-/Fockrechnung getrennt; der gemeinsame Replay führt beide
normal und optimiert aus.

\subsection{15. Fortgesetzte Auswahlprüfung: fehlt eine
Spiegelungsrelation?}\label{fortgesetzte-auswahlpruxfcfung-fehlt-eine-spiegelungsrelation}

Der Nutzer bat während der Konsolidierung um weitere Arbeit. Deshalb
wurde nicht nur die neue Quelle beschrieben, sondern unmittelbar die
wichtigste noch verdeckte Wahl angegriffen: \textbf{Warum gerade die
gleichstarke Nachbarmischung?}

\subsubsection{15.1 Ein Gegenmodell gegen eine zu starke Deutung des
Grundzustandssatzes}\label{ein-gegenmodell-gegen-eine-zu-starke-deutung-des-grundzustandssatzes}

Ersetze die festgelegte Quelle zunächst durch die einfache reelle
Familie

\[\adjustbox{max width=\linewidth}{$\displaystyle U_\eta(a,b)=aI+bR_\eta,\qquad a^2+b^2=1.$}\]

Alle einzelnen Quellen bleiben CAR-normiert, und \([U_\eta,R_\eta]=0\)
gilt für jedes a und b. Der gemeinsame Clock und die volle innere
Symmetrie bleiben erhalten. Ihre Einteilchenspektren sind

\[\adjustbox{max width=\linewidth}{$\displaystyle \operatorname{spec}S_+=\{1-2ab,1,1,1+2ab\},$}\]

\[\adjustbox{max width=\linewidth}{$\displaystyle \operatorname{spec}S_-=\{1-\sqrt2ab,1-\sqrt2ab,1+\sqrt2ab,1+\sqrt2ab\}.$}\]

Für den periodischen Lift und \((a,b)=(3/5,4/5)\) ist der kleinste
Eigenwert \(1/25>0\). Es gibt keine freie Fermionrichtung, obwohl das
Umlaufvorzeichen positiv ist. Der Grundzustandsbeweis aus Abschnitt 6
gilt unverändert mit \(\delta=(15/2)/25=3/10\). Bei
\(0<|g|/\Delta\le1/10000\) ist seine Ladungstrennung mindestens

\[\adjustbox{max width=\linewidth}{$\displaystyle \left(\frac3{10}-\frac{1920^2}{10000^2}\right)\frac{g^2}{\Delta}
=\frac{8223}{31250}\frac{g^2}{\Delta}>0.$}\]

Damit hat \textbf{auch dieser unverdrehte, ungleich gewichtete Kandidat}
einen eindeutigen vollen Grundzustand bei N=256. Die Paarmatrix ist hier
\(8I+8a^2b^2 A_{C_4}\) und nicht dieselbe wie bei der zuvor fixierten
gleichstarken Quelle. Das ist also kein Gegenbeispiel gegen die Sätze
der Abschnitte 3-7, aber eines gegen die weitergehende Behauptung „Clock
plus eindeutiger Grundzustand erzwingen den negativen Lift``. Die
Gewichtung muss mit ausgewählt werden.

\subsubsection{15.2 Eine kleine Spiegelungsregel würde die Gewichtung
bestimmen}\label{eine-kleine-spiegelungsregel-wuxfcrde-die-gewichtung-bestimmen}

Für beide Lifts wähle den reellen signierten Reflexionsoperator J durch

\[\adjustbox{max width=\linewidth}{$\displaystyle Jf_0=f_0,\qquad Jf_x=\eta f_{4-x}\quad(x=1,2,3).$}\]

Er erfüllt exakt \(J^2=I\) und \(JR_\eta J=R_\eta^{-1}\). Setze
\(L=R_\eta J\); auch \(L^2=I\). Auf dem markierten Zyklus entspricht J
einer Vertexreflexion, L der dazu verschobenen Kantenreflexion. Die
zugehörigen Paar-/Bosonbänke werden mit der vorzeichenlosen Permutation
von L reflektiert. Dies ist der ausdrücklich erklärte gemeinsame
Reflexionsvertrag, nicht irgendeine frei austauschbare Basis.

Nun gilt die kurze Identität

\[\adjustbox{max width=\linewidth}{$\displaystyle \boxed{U(a,b)J=L\,U(b,a).}$}\]

Die Spiegelung vertauscht also genau die beiden Quellgewichte. Beim
ursprünglichen Paarvertex sind die Kopienkoeffizienten die äußeren
Produkte der Quellenzeilen mit sich selbst, ohne komplexe Konjugation.
Die Gleichheit der reflektierten Paaroperatoren verlangt deshalb

\[\adjustbox{max width=\linewidth}{$\displaystyle a^2=b^2.$}\]

Notwendigkeit sieht man bereits an den Koeffizienten, die zwei interne
Fermionen vom linken beziehungsweise rechten Kopienort entnehmen. Für
Suffizienz gilt bei \(a=b\) bereits \(UJ=LU\), bei \(a=-b\) entsprechend
\(UJ=-LU\); das gemeinsame Minus verschwindet im Paar. Zusammen mit der
Normierung folgt \(|a|=|b|=1/\sqrt2\). Die vier Zeilen und ihre
Paarprodukte wurden exakt für beide Lifts geprüft.

\textbf{Bedingter Auswahlsatz:} Innerhalb der reellen, gleichförmigen
Zweinachbar-Klasse wählen (i) die erklärte gemeinsame
Vertex-/Kantenreflexion und (ii) keine exakt freie Fermionrichtung die
balancierte antiperiodische Quelle aus, bis auf Quellenzeichen und
entsprechende lokale Vorzeichenkonventionen. Denn die Reflexion erzwingt
Balance; bei Balance hat nur der negative Viererlift vollen Rang. Für
vier Knoten sind die beiden relativen Vorzeichen der balancierten
Mischung durch alternierende lokale Fermionvorzeichen äquivalent. Die
zugehörigen Bosonzeichen dürfen dabei unverändert bleiben, weil der
Vertex gerade ist.

Das ist eine kleinere fehlende Bedingung als eine beliebige neue
Dynamik: \textbf{Wende dieselbe Spiegelung auf Fermionorte und
Paarübergänge konsistent an.} Sie ist dennoch eine Zusatzvoraussetzung,
solange sie nicht aus der eigentlichen Quelle stammt. Das Fehlen freier
Zuschauer ist ebenfalls ein hier erklärtes Auswahlkriterium, nicht ein
allgemeines Naturgesetz.

\subsubsection{15.3 Tatsächlicher Anschluss an ältere
TFPT-Struktur}\label{tatsuxe4chlicher-anschluss-an-uxe4ltere-tfpt-struktur}

In \codeword{origin\_theory.tex} ist die Relation
\(\sigma\rho\sigma=\rho^{-1}\) bereits Teil der geometrischen
Normalform. \codeword{v177\_seam\_marking\_kernel.py} verwendet explizit
\(\rho:z\mapsto iz\) und \(\sigma:z\mapsto1/z\) auf den vier Marken.
\codeword{v180\_clock\_is\_mobius.py} trennt die Clock-Geometrie
weiterhin von der ursprünglichen Realisierungsvoraussetzung. Diese
beiden Dateien wurden für diesen Anschluss gelesen, nicht als
vollständiger neuer Seam-Beweis ausgeführt.

Somit wird die Spiegelung nicht als neue Zahlenidee importiert. Es fehlt
aber die entscheidende \textbf{Operatorzuordnung}: Die geometrische
sigma muss auf den ursprünglichen fermionischen Ressourcen als das
passende J wirken und zugleich die W-Paarbänke mit der verschobenen
Kantenwirkung von L abbilden. Die abstrakte dihedrale Relation allein
bestimmt diese relative Wirkung auf zwei verschiedenen Ressourcen nicht.
Auch Reflexionspositivität ist nicht dasselbe wie diese unitäre
Paarvertex-Symmetrie. Ein physischer Lorentz-/Pin-Adapter wird aus den
reellen Vorzeichenmatrizen nicht behauptet.

Eine zusätzliche exakte Negativkontrolle macht diese Grenze sichtbar:
Reflektiert man die Paarbänke stattdessen mit \textbf{derselben
Vertexpermutation} wie die Fermionorte, verlangt die Paar-Kovarianz in
diesem Ansatz \(b=0\). Schon ein diagonaler Kopienkoeffizient des
Differenzoperators ist \(b^2\). Diese andere Zuordnung erlaubt also nur
die entkoppelte Vor-Ort-Quelle statt der balancierten Nachbarquelle. Die
Wahl zwischen Vertex- und Kantenwirkung ist physischer Inhalt, nicht
durch das Wort „Spiegelung`` erledigt.

Der nächste besonders aussagekräftige Test lautet damit enger:
\textbf{Liefert die ursprüngliche markierte TFPT-Seam genau dieses
gemeinsame Paar von Spiegelungswirkungen, einschließlich Zustand und
Vertex?} Ein positiver Nachweis würde in dieser Quellenklasse die
Mischungswahl und das Umlaufvorzeichen zusammen ersetzen. Ein negativer
Nachweis würde diesen konkreten Anschluss ausschließen. Der
Kopplungswert, die Verfeinerungsregel und die T1-T8-Gesamtnachweise
blieben davon getrennte Aufgaben.

\subsubsection{Abgeschlossener eigener
Replay}\label{abgeschlossener-eigener-replay}

1669 explizite Bedingungen normal und optimiert, byteidentische
Ergebnisdateien. Darin enthalten sind wiederholte Quellen-, Vorzeichen-
und CAR-Prüfungen; die Zahl zählt keine unabhängigen Theoreme. Der
allgemeine Grundzustands- und Antwortbeweis ist in Abschnitten 5-7
ausgeschrieben. Die Zusatzprüfung der beiden späten Perspektivtexte
umfasst 61 weitere Bedingungen, ebenfalls normal/optimiert
byteidentisch. Ihre allgemeinen Voraussetzungen und Grenzen stehen in
Abschnitt 14. Die anschließende Auswahlprüfung umfasst 45 weitere exakte
Bedingungen; Gegenmodell und bedingter Spiegelungs-Auswahlsatz stehen in
Abschnitt 15.

\newpage

